\documentclass[10pt]{article}
\usepackage[margin=20mm]{geometry}
\usepackage{amsmath,amssymb,fontspec,unicode-math,xurl,hyperref}
\IfFontExistsTF{DejaVu Serif}{\setmainfont{DejaVu Serif}}{\setmainfont{Latin Modern Roman}}
\setmathfont{latinmodern-math.otf}
\setlength{\parindent}{0pt}\setlength{\parskip}{6pt}\sloppy
\title{Technical Verification of Behavioral Exploration: Learning to Explore via In-Context Adaptation}
\author{AI-assisted technical review}
\date{}
\begin{document}
\maketitle
Abstract: All 28 physical pages, including the proof and experimental appendix, were read. The finite-support argument establishes balanced cumulative visitation, but its conversion to a single-episode optimal policy requires an explicit mixture. The experimental appendix also contains an encoder inconsistency and does not specify the numerical deployment coverage conditioning. These are Minor issues. The supplied evidence supports the reported qualitative advantage in the stated exploration settings; it does not independently reproduce the numerical results. Scientific adjudication remains pending.\par
\section*{Identity and coverage}
This report concerns Andrew Wagenmaker, Zhiyuan Zhou and Sergey Levine, ICML 2025, PMLR 267, I25-43764 (0395), OpenReview `tlLkY9E2bZ`. Citations use one-based physical pages, which agree with the printed numbers. The original PDF and its embedded Appendices A and B were inspected. All numbered Figures 1–31 were actually viewed, including the per-task grids; these grids were checked for axes, monotonicity and visible exceptions, not digitized into new measurements. Tables 1–9 were read, and the architecture discrepancy was checked in source pixels.\par
The frozen source manifest reports no identified public implementation after bounded discovery; this is scoped non-discovery, not a claim that no implementation exists. The separate OpenReview attachment was inaccessible in the acquisition record. No author program, robot, simulator, model or dataset was executed or downloaded. The only executable check is a reviewer-written finite example and table arithmetic.\par
\section*{Supported findings}
Problem 1 (The optimal output policy needs an explicit episode mixture).  Objective 4.1 (p. 4) optimizes \(\operatorname{cov}_{\beta}(\Lambda_\pi)\) for the expected feature matrix of one trajectory. Proposition A.1 (pp. 15–16) instead proves that an adaptive sequence of \(n_\beta\) episodes visits each supported terminal feature once, then concludes that the terminal-feature probabilities are \(1/n_\beta\). The latter does not follow for the last episode, nor for each individual episode, without specifying how the output episode is selected.\par
Consider a deterministic one-step MDP with two actions leading to terminal features \(e_1,e_2\), zero nonterminal features, a fixed initial state, and behavior probabilities \((0.9,0.1)\). This satisfies the formal assumptions with \(n_\beta=2\), \(0<\epsilon<0.1\), and finite actions. Initially either terminal gives maximal coverage, so the conditioned first-episode distribution is still \((0.9,0.1)\). In episode two the policy chooses the previously unvisited terminal, giving marginal probabilities \((0.1,0.9)\), not \((1/2,1/2)\). For \(\lambda=0.01\), the one-episode objective is \[f(p)=\left((p+\lambda)^{-1}+(1-p+\lambda)^{-1}\right)^{-1},
 \qquad f(0.9)=0.098137\ldots< f(1/2)=0.255.\] This refutes the individual-episode interpretation of the final probability statement. It does not refute the cumulative-coverage induction: both features are visited exactly once. A uniformly sampled episode from the two-episode interaction has the optimal distribution. More generally, specify the output as the uniform mixture of the \(n_\beta\) episode policies, including their induced histories, or define the objective using normalized cumulative occupancy. Then the proof’s conclusion follows by linearity and the convexity of \(x\mapsto1/(x+\lambda)\). The phrase “after playing for \(n_\beta\) episodes” plausibly intends precisely this interpretation, so the established finding is a Minor output-policy specification ambiguity. The example does not refute the intended cumulative-coverage guarantee, and the original per-episode optimality interpretation is not certified. The main informal proposition also replaces zero nonterminal features by a common extra coordinate; its fixed contribution does not disturb the balancing argument.\par
The internal induction was checked under its stated deterministic, finite-action, positive-support assumptions: before the last round an unvisited supported terminal exists, its conditioning event has positive probability, and an action with no supported continuation to such a terminal has conditional mass zero. The notation \(w^{\pi_\beta}(s)\) is defined for terminal visitation but later applied to intermediate states; a time-indexed reachability probability would avoid this additional local notational defect.\par
Problem 2 (The deployed policy is not fully specified by the appendix).  Section B.2 (p. 20) says the image encoder is ResNet-18, whereas Table 7 (p. 21) says ResNet-34 for both BE and BC. In addition, the deployment rule on p. 5 requires a scalar \(\mathrm{exp}\), but the full experimental appendix does not state its numerical value or its normalization for each setting. Figure 6 (p. 7) labels this value as a percentage without defining its mapping to the absolute inverse-trace coverage.\par
These are concrete reconstruction gaps. The regularizer \(\lambda=0.01\) and 32-dimensional random feature network are supplied for D4RL (p. 18); Libero instead uses 16 cosine features (p. 20), and WidowX uses relative coordinates (p. 26). Consequently a percentage cannot simply be assumed to have a universal absolute meaning. Resolve the encoder discrepancy and supply the actual configuration, coverage preprocessing, chosen deployment target, regularization for the vision settings, training duration and checkpoint-selection rule. This Minor finding concerns reproducibility; no available implementation resolves it, and no causal effect on the reported ordering has been established. Table 2’s caption also incorrectly refers to Figures 7–8 although its rows concern D4RL Figures 3–5.\par
\section*{Checks and limits of the evidence}
The theory depends on exact behavior-conditional probabilities and deterministic, terminal-only one-hot features. Equation (4)’s finite-data diffusion approximation, random continuous features, and deployment history sampling are empirical extensions. Section 4.3 explicitly acknowledges the history-distribution fixed point and adopts uniformly sampled demonstration histories. The proof is therefore not a guarantee for all robotic experiments; the paper does not supply such a guarantee.\par
The D4RL evaluation removes offline rewards, shortens Antmaze episodes, and measures first accomplishment or region coverage rather than standard asymptotic return. These changes are disclosed. Figures 3–5 and Table 2 agree that BE leads on aggregated Antmaze goals and regions, while SUPE has the higher final Kitchen value (0.956 versus 0.922); the text appropriately says comparable performance there. The detailed Antmaze panels reveal exceptions, including a medium-maze goal on which BC leads, so the support is for the aggregate comparison. The two-terminal calculation gives BC an expected \(1+0.9/0.1+0.1/0.9=10.111\ldots\) episodes to see both terminals, versus two for ideal BE, verifying the motivating distinction without simulating the paper’s systems.\par
Libero uses the same 90 demonstrated tasks for evaluation, with hidden or provided task labels, 15 or six attempts, respectively. This tests adaptive selection among known behaviors, not held-out-task generalization. Table 6 gives a hidden-task improvement of 23.4 percentage points over BC, versus only 1.6 points in the task-provided setting. BC action noise and BC history controls are included. However, Table 7 uses action chunks of 20 for BE and 10 for BC, and tuning chooses best-performing architectures. These disclosed implementation differences limit attribution to the coverage variable alone. A matched-chunk ablation would strengthen the mechanism claim; it is not evidence that the whole advantage disappears. Error bars need the aggregation unit made explicit: five training seeds, 90 tasks and three rollouts are clustered. The reported 0.010 standard error at success 0.824 is numerically compatible with pooling 1350 Bernoulli observations (0.01036), but this coincidence does not establish the actual procedure. Seed-level or hierarchical uncertainty requires raw per-seed results.\par
WidowX has three scenes and eight trials per scene. A success means touching both objects within five attempts; a failure is explicitly scored as six attempts for the capped attempt metric. Thus the figure is not an uncensored time-to-success estimate. Source bars in Figures 27–29 are consistent with success counts BE \((5,8,6)\), BC \((0,7,0)\) and OpenVLA \((3,2,5)\) out of eight, giving aggregate percentages 79.17, 29.17 and 41.67, consistent with Figure 9. These are inferred from plotted increments, not recovered raw trials. The text’s “roughly 40\% more” is better expressed in percentage points and by comparator. The OpenVLA prompts and additional “other silver pot” probe are disclosed; this comparison is specific to ambiguous prompts and does not establish general inferiority on manipulation tasks. Figures 30–31 visibly support the qualitative history-adaptation illustration, without independently verifying training or rollout provenance.\par
\section*{Repair and verification status}
The necessary mathematical repair is to name the cumulative or mixture objective and state the output-policy construction. The configuration repair requires the authors’ actual experiment settings; this review has not guessed or substituted them. No Major defect was established within the inspected evidence. The two Minor findings and the explicit uncertainty requests remain pending. A fresh independent acceptance check revisited both Problems, including the cumulative interpretation and the boundaries of the finite example. No human review is recorded.\par
\end{document}
